\documentclass[11pt,a4paper]{article}
\usepackage[margin=0.8in,top=0.85in,bottom=0.85in]{geometry}
\usepackage{amsmath,amssymb,amsfonts}
\usepackage{bm}
\usepackage{booktabs}
\usepackage{enumitem}
\usepackage{fancyhdr}
\usepackage{microtype}
\usepackage{xcolor}
\usepackage[colorlinks=true,linkcolor=black,citecolor=black,urlcolor=black]{hyperref}

% Clean, airy typography with generous breathing room
\linespread{1.06}
\setlength{\parindent}{0pt}
\setlength{\parskip}{4.5pt plus 1pt minus 1pt}

% Header and Footer
\pagestyle{fancy}
\fancyhf{}
\renewcommand{\headrulewidth}{0pt}
\rfoot{\footnotesize Page \thepage\ of 2}
\lfoot{\footnotesize \textit{IIT Bombay --- Department of Electrical Engineering --- EE 603: Digital Signal Processing}}

\begin{document}

% ═════════════════════════════════════════════════════════════════
% IIT BOMBAY EE 603 DEPARTMENTAL HEADER
% ═════════════════════════════════════════════════════════════════
\noindent
\begin{tabular*}{\textwidth}{@{}l @{\extracolsep{\fill}} r@{}}
\textbf{\large INDIAN INSTITUTE OF TECHNOLOGY BOMBAY} & \textbf{Autumn 2026} \\
Department of Electrical Engineering & Filter Design Laboratory \\
\textbf{EE 603: Digital Signal Processing} & Wiki \& Contributor Guide
\end{tabular*}

\vspace{6pt}
\hrule height 1.4pt \vspace{2pt} \hrule height 0.5pt
\vspace{20pt}

\begin{center}
{\LARGE \textbf{Overtune 3: Architecture, Wiki \& Contributor Guide}}\\[8pt]
{\large \textsc{DSP Engine Internals, Numerical Stability Guidelines, and Development Protocol}}
\end{center}

\vspace{18pt}

\section*{1. Project Architecture \& Repository Layout}

Overtune~3 is developed with a strict decoupling between the mathematical DSP synthesis engine (pure C++20) and the desktop presentation layer (Qt~6 QML and Modern C++ wrapper classes).

\vspace{4pt}
\begin{verbatim}
overtune3/
|-- dsp/                     <-- Pure C++20 signal processing library (no GUI dependencies)
|   |-- include/overtune/    <-- Public headers (FilterDesigner, Biquad, ComplexPoly)
|   +-- src/                 <-- Mathematical implementations and code generators
|-- app/                     <-- Desktop GUI application (Qt 6.4+, QML, Quick Controls)
|   |-- src/                 <-- C++ QObject models (FilterEngine, ExportModel, ThemeManager)
|   +-- qml/                 <-- Declarative UI (Pages, Controls, Custom Canvas Plots)
|-- latex/                   <-- Authoritative documentation sources and compiled PDFs
+-- tests/                   <-- Numerical test suites verified against SciPy benchmarks
\end{verbatim}

\vspace{6pt}
\section*{2. The Core DSP Engine (\texttt{dsp/})}

The synthesis pipeline is encapsulated within the \texttt{overtune::FilterDesigner} class:
\begin{itemize}[leftmargin=20pt,itemsep=2pt]
\item \textbf{Prototype Synthesizer:} Generates analog prototype poles and zeros for normalized low-pass prototypes in the continuous $s$-domain.
\item \textbf{Frequency Transformation:} Performs LP-to-HP, LP-to-BP, or LP-to-BS conformal mapping.
\item \textbf{Bilinear Transform ($s \to z$):} Maps poles and zeros to discrete coordinates using pre-warped angular frequencies $\Omega = 2 f_s \tan(\omega / 2)$.
\item \textbf{SOS Pairing \& Ordering:} Decomposes the high-order transfer function into Second-Order Sections (SOS). Complex conjugate pole pairs with the highest $Q$ are paired with nearest transmission zeros and placed downstream to maximize dynamic range and minimize register overflow.
\end{itemize}

\section*{3. Implementing a New Filter Family}

To contribute a new filter approximation (e.g., Papoulis 'L' filter, Legendre, or Halpern):
\begin{enumerate}[leftmargin=20pt,itemsep=2pt]
\item Define the analog prototype root-finding algorithm in \texttt{dsp/src/Prototypes.cpp}.
\item Return a \texttt{std::vector<std::complex<double>>} of normalized $s$-plane poles and zeros.
\item Register the prototype enum identifier in \texttt{dsp/include/overtune/Types.h}.
\item Expose the property binding in \texttt{app/src/FilterEngine.h} and add the corresponding selector option in \texttt{app/qml/pages/DesignPage.qml}.
\item Write a unit test in \texttt{tests/test\_prototypes.cpp} verifying magnitude response error $\le 10^{-5}$ relative to SciPy reference vectors.
\end{enumerate}

\section*{4. Numerical Accuracy \& Stability Standards}

All contributions must adhere to the following numerical standards:
\begin{itemize}[leftmargin=20pt,itemsep=2pt]
\item \textbf{Double-Precision Computation:} All internal prototype polynomials, matrix inversions, and root-finding iterations must use IEEE 754 64-bit double precision (\texttt{double}).
\item \textbf{SOS Decomposition:} High-order filters must never be exported as single high-degree rational transfer functions $B(z)/A(z)$ due to root ill-conditioning for $N > 6$.
\item \textbf{Jury Stability Verification:} Every synthesized filter must pass $|p_k| < 1 - 10^{-7}$ for all poles before being emitted to the UI or code exporter.
\item \textbf{Zero Heap Allocation in Audio Thread:} Processing callbacks in \texttt{processBlock()} must execute in strictly bounded deterministic time with zero memory allocations or system calls.
\end{itemize}

\section*{5. Canonical DSP Literature References}
\begin{table}[h!]
\centering
\small
\renewcommand{\arraystretch}{1.15}
\begin{tabular}{lll}
\toprule
\textbf{Author(s)} & \textbf{Title} & \textbf{Core Topic} \\
\midrule
A.~V.~Oppenheim \& R.~W.~Schafer & \textit{Discrete-Time Signal Processing} & Canonical LTI theory, Bilinear transform \\
S.~K.~Mitra & \textit{Digital Signal Processing} & Comprehensive filter realization structures \\
L.~R.~Rabiner \& B.~Gold & \textit{Theory and Application of DSP} & Filter approximation algorithms \\
A.~Antoniou & \textit{Digital Signal Processing} & Elliptic rational filter synthesis \\
R.~Bristow-Johnson & \textit{Audio EQ Cookbook} & Biquad parametric EQ derivations \\
\bottomrule
\end{tabular}
\caption{Foundational reference literature for Overtune 3 developers.}
\end{table}

\section*{6. Git Workflow \& Contribution Protocol}
\begin{itemize}[leftmargin=20pt,itemsep=2pt]
\item \textbf{Branching Model:} Create feature branches branching from \texttt{main} named \texttt{feature/<filter-name>} or \texttt{fix/<issue-number>}.
\item \textbf{Formatting:} Format all C++ code with \texttt{clang-format -style=file} adhering to the repository `.clang-format`.
\item \textbf{Continuous Integration:} Automated GitHub Actions build on Linux, macOS, and Windows with clang and gcc. All regression unit tests must pass.
\item \textbf{Pull Requests:} PRs must include a clear mathematical description of any proposed filter topology, accompanied by comparison plots generated from the test suite.
\end{itemize}

\end{document}
